% =============================================================================
\section{Dataset Privacy Risk Index}
\label{sec:dpri}
% =============================================================================

Motivated by Theorem~\ref{thm:main}, we define five distributional
features that are computable from raw data before any model training.
Each feature targets a geometric property that the theorem identifies as
a driver of MIA advantage.

\subsection{Feature Definitions}

Let $D = \{x_1,\dots,x_n\}$ denote the feature matrix (labels used only
for Feature~5).
All per-sample scores are aggregated to dataset level using mean and
90th percentile, yielding a feature vector
$\phi(D) \in \mathbb{R}^{8}$ (two statistics for each of the four
per-sample features, plus two scalars).

\paragraph{Feature 1: Sample Uniqueness $u(x)$.}
\begin{equation}
  u(x) = \min_{x' \in D,\, x'\neq x} \|x - x'\|_2
\end{equation}
High $u(x)$ means $x$ is isolated from all other training points.
Isolated samples have high MIA advantage by Equation~\eqref{eq:bound}.
We use the 5-NN distance as a robust estimate. The nearest neighbour
alone ($k=1$) is noisy for dense datasets.

\paragraph{Feature 2: Local Density $\rho(x)$.}
\begin{equation}
  \rho(x) = \frac{1}{r_k(x)}, \quad
  r_k(x) = \text{distance to } k\text{-th nearest neighbour}
\end{equation}
Low $\rho(x)$ (sparse region) implies large $u(x)$ and high risk.
The two features are theoretically related but empirically complementary:
$u(x)$ captures global isolation while $\rho(x)$ captures local structure.

\paragraph{Feature 3: Outlier Score $o(x)$.}
\begin{equation}
  o(x) = \frac{1}{2}\bigl(\mathrm{LOF}_{\mathrm{norm}}(x)
                         + \mathrm{iForest}_{\mathrm{norm}}(x)\bigr)
\end{equation}
An ensemble of Local Outlier Factor~\cite{breunig2000lof} and Isolation
Forest~\cite{liu2008isolation}, each normalized to $[0,1]$.
The ensemble reduces the sensitivity of any single outlier detector to
its specific density assumption.

\paragraph{Feature 4: Feature Entropy $H(X)$.}
\begin{equation}
  H(X) = \frac{1}{d} \sum_{j=1}^{d} H(X_j), \quad
  H(X_j) = -\sum_b p_b \log_2 p_b
\end{equation}
The average Shannon entropy across all features $X_j$, computed over
a histogram discretization.
High-entropy features increase the dimensionality of the space in which
samples are embedded, reducing density and increasing uniqueness.

\paragraph{Feature 5: Cluster Separation $S$.}
\begin{equation}
  S = \mathrm{Silhouette}(X, y)
  = \frac{1}{n}\sum_{i=1}^{n} \frac{b(x_i) - a(x_i)}{\max(a(x_i), b(x_i))}
\end{equation}
where $a(x_i)$ is the mean intra-class distance and $b(x_i)$ is the
mean nearest-class distance.
Well-separated classes ($S$ near $1$) place decision boundary samples
in high-contrast regions, increasing their MIA risk.

\subsection{Aggregation and the DPRI Score}

The per-sample scores $(u, \rho, o)$ are aggregated to dataset level as
mean and 90th percentile, giving a feature vector:
\begin{equation}
  \phi(D) = \bigl(
    \bar{u},\; u_{90},\;
    \bar{\rho},\; \rho_{90},\;
    \bar{o},\; o_{90},\;
    H(X),\; S
  \bigr) \in \mathbb{R}^{8}.
\end{equation}
The 90th percentile captures tail behaviour (the hardest-to-protect
samples) which disproportionately drives aggregate MIA AUC.

The final DPRI score is a linear combination:
\begin{equation}
  \mathrm{DPRI}(D) = w^\top \phi(D),
  \label{eq:dpri}
\end{equation}
where $w \in \mathbb{R}^{8}$ is learned by ridge regression to predict
$\mathrm{Risk}(D)$ (Section~\ref{sec:setup}).
The linear form is intentional: it is interpretable, and
Section~\ref{sec:results} shows it performs comparably to nonlinear
regressors, indicating the relationship between DPRI features and
$\mathrm{Risk}(D)$ is approximately linear.

\subsection{Computation Cost}

All five features require only standard nearest-neighbour and
density-estimation routines.
On a laptop (8-core CPU), DPRI computation takes \TODO{X} seconds for
datasets up to 50k samples and \TODO{X} minutes for 200k samples.
No GPU is required.
This is the key practical advantage over MIA-based risk assessment,
which requires training and attacking multiple models.
